% Part: computability
% Chapter: recursive-functions
% Section: notations-pr-functions

\documentclass[../../../include/open-logic-section]{subfiles}

\begin{document}

\olfileid{cmp}{rec}{not}
\olsection{আদিম পুনরাবৃত্ত অপেক্ষকের সংকেতলিপি}

আদিম পুনরাবৃত্ত অপেক্ষকগুলির সুনির্দিষ্ট আবেশনভিত্তিক বর্ণনা থাকার একটি
সুবিধা হলো, আমরা সুশৃঙ্খলভাবে তাদের বর্ণনা করতে পারি। যেমন, প্রত্যেকটি
এমন অপেক্ষককে নিচের মতো একটি “সংকেতলিপি” দেওয়া যায়। শূন্য, উত্তরসূরি
ও অভিক্ষেপগুলির জন্য যথাক্রমে $\Zero$, $\Succ$ ও $\Proj{n}{i}$ প্রতীক
ব্যবহার করি। এবার ধরা যাক, একটি $k$-স্থানীয় অপেক্ষক~$f$ ও
$n$-স্থানীয় অপেক্ষক $g_0$, \dots,~$g_{k-1}$ থেকে মিশ্রণের সাহায্যে
$h$ সংজ্ঞায়িত, এবং শেষোক্ত অপেক্ষকগুলিকে আমরা যথাক্রমে $F$, $G_0$,
\dots,~$G_{k-1}$ সংকেতলিপি দিয়েছি। তাহলে নতুন প্রতীক
$\fn{Comp}_{k,n}$ ব্যবহার করে $h$ অপেক্ষকটিকে
$\fn{Comp}_{k,n}[F,G_0,\dots,G_{k-1}]$ দিয়ে বোঝানো যায়।

আদিম পুনরাবৃত্তির মাধ্যমে সংজ্ঞায়িত অপেক্ষকগুলির জন্য অনুরূপ সংকেতলিপি
ব্যবহার করা যায়। ধরা যাক, $(k+1)$-স্থানীয় অপেক্ষক~$h$-কে
$k$-স্থানীয় অপেক্ষক~$f$ ও $(k+2)$-স্থানীয় অপেক্ষক~$g$ থেকে আদিম
পুনরাবৃত্তির মাধ্যমে সংজ্ঞায়িত করা হয়েছে, এবং $f$ ও~$g$-কে দেওয়া
সংকেতলিপি যথাক্রমে $F$ ও~$G$। তাহলে~$h$-কে দেওয়া সংকেতলিপি হলো
$\fn{Rec}_k[F,G]$।

মনে করুন, যোগের অপেক্ষকটি আদিম পুনরাবৃত্তির সাহায্যে এভাবে সংজ্ঞায়িত:
\begin{align*}
  \Add(x_0, 0) & = \Proj{1}{0}(x_0) = x_0\\
  \Add(x_0, y+1) & = \Succ(\Proj{3}{2}(x_0, y, \Add(x_0, y))) = \Add(x_0, y) +1
\end{align*}
এখানে~$f$-এর ভূমিকা পালন করে $\Proj{1}{0}$, আর~$g$-এর ভূমিকা পালন করে
$\Succ(\Proj{3}{2}(x_0, y, z))$। শেষের অপেক্ষকটিকে
$\fn{Comp}_{1,3}[\Succ,\Proj{3}{2}]$ সংকেতলিপি দেওয়া হয়, কারণ এটি
$1$-স্থানীয় অপেক্ষক~$\Succ$ ও $3$-স্থানীয় অপেক্ষক~$\Proj{3}{2}$ থেকে
মিশ্রণের সাহায্যে একটি অপেক্ষক সংজ্ঞায়িত করার ফল। এই ব্যবস্থা অনুযায়ী
যোগের অপেক্ষকটিকে আমরা এভাবে বোঝাতে পারি:
\[
\fn{Rec}_1[\Proj{1}{0},\fn{Comp}_{1,3}[\Succ,\Proj{3}{2}]].
\]
এই সংকেতলিপিগুলি কখনো কখনো কাজে লাগে; যেমন, আদিম পুনরাবৃত্ত অপেক্ষকগুলি
গণনা করে তালিকাভুক্ত করার সময়।

\begin{prob}
  $\Mult$-এর সম্পূর্ণ আদিম পুনরাবৃত্ত সংকেতলিপি দিন।
\end{prob}

\end{document}
